%% toc: P3 40 40.2
\chapter{Random Graphs and Extension Statements}
\label{sec:p03-w03-p03-065-random-graphs}

The zero--one law of \cref{sec:p03-zero-one-law} concerns one probability
space for each size \(n\): the uniform labelled graphs, in which every pair of
vertices is an edge with probability \(1/2\).  Random graph theory lets the
edge probability depend on \(n\).  As the probability grows, a property passes
from almost never holding to almost surely holding, and the place where this
happens is its threshold.  This chapter sets up the random graph
\(G(n,p(n))\), proves the threshold behaviour of the first-order examples
that motivate the subject, and then studies the first-order sentences that
carry the probabilistic analysis of the sparse regime \(p=n^{-\alpha}\): the
extension statements.  The main result is the criterion that, for irrational
\(\alpha\in(0,1)\), an extension statement holds almost surely exactly when its
rooted graph is safe.

\paragraph{Inputs and conclusions.}
The input is an edge-probability function \(p=p(n)\) with values in
\([0,1]\) and a first-order property of graphs, given either as a sentence or
as a rooted graph \((R,H)\) together with an irrational exponent
\(\alpha\in(0,1)\).  The conclusions are statements about
\(\lim_{n\to\infty}\Pr[G(n,p(n))\entails A]\): that it is \(0\), that it is
\(1\), or where, as \(p(n)\) varies, it changes from one to the other.  Each
conclusion is asymptotic.  None of them decides whether a particular finite
graph has the property, and none of them decides the property for any
particular \(n\).

\paragraph{Selected treatment.}
The definitions, examples and source claims follow Joel Spencer's chapter
\emph{Logic and Random Structures}, in the 30-page author's version hosted on
Spencer's New York University page
\cite{Spencer2007LogicRandomStructuresAuthor}.  The printed chapter in the
2007 Springer volume \emph{Finite Model Theory and Its Applications} has not
been consulted.  Every locator below refers to the author's version, whose
sections are numbered \(1\) to \(14\), followed by an unnumbered section on
the references, and whose pages are numbered \(1\) to \(30\).  The random
graph and its threshold examples are its Section~2
\cite[pp.~4--6]{Spencer2007LogicRandomStructuresAuthor}; extension statements
are its Section~3 \cite[pp.~6--8]{Spencer2007LogicRandomStructuresAuthor};
the counting and generic extension statements behind its almost sure theory
are in its Section~5 \cite[pp.~9--11]{Spencer2007LogicRandomStructuresAuthor}.
The author's version gives heuristic justifications for the thresholds,
proves only one direction of the extension criterion, saying the other needs
more work \cite[p.~7]{Spencer2007LogicRandomStructuresAuthor}, and omits the
proof of its large deviation estimate
\cite[pp.~10--11]{Spencer2007LogicRandomStructuresAuthor}.  The proofs in the
annex are complete proofs written for this chapter.  The extension criterion
is proved there by a correlation inequality of Janson type
(\cref{thm:p03-w03-p03-065-janson}), itself derived from the Harris
inequality (\cref{thm:p03-w03-p03-065-harris}); the counting theorem is
proved by high centred moments.  Two slips in the worked examples of the
author's version are corrected in \cref{def:p03-w03-p03-065-pi-seven}.

\section{The random graph and threshold functions}
\label{sec:p03-w03-p03-065-thresholds}

\begin{hbdefinition}{def:p03-w03-p03-065-gnp}{The random graph \(G(n,p)\)}
A \emph{graph} is a set \(V\) of vertices with an irreflexive symmetric
binary relation \(\sim\) on \(V\), the adjacency relation; \(x\sim y\) is read
``\(x\) and \(y\) are adjacent''.  For an integer \(n\geq1\) and a real
\(p\in[0,1]\), the random graph \(G(n,p)\) has vertex set \(V=[n]\), and for
distinct \(x,y\) the events \(x\sim y\) have probability \(p\) and are
mutually independent \cite[p.~4]{Spencer2007LogicRandomStructuresAuthor}.  Formally the
probability space is the set of all graphs on \([n]\), and a graph with
exactly \(m\) edges has probability \(p^m(1-p)^{\binom n2-m}\).  The graphs
are undirected, without loops and without multiple edges.

Throughout this chapter \(p=p(n)\) is a function of \(n\), and \(G(n,p(n))\)
is the sequence of these probability spaces.  For \(p=1/2\) the space is the
uniform space on labelled graphs of \cref{sec:p03-zero-one-law}.
\end{hbdefinition}

\begin{hbdefinition}{def:p03-w03-p03-065-almost-sure}{Almost surely, almost never}
Let \(A\) be a property of graphs, for example the class of models of a
first-order sentence in the language \(\{\sim\}\) with equality.  With
respect to a given \(p(n)\), \(A\) holds \emph{almost surely} if
\(\lim_{n\to\infty}\Pr[G(n,p(n))\entails A]=1\), and \(A\) holds
\emph{almost never} if this limit is \(0\), equivalently if \(\neg A\) holds
almost surely \cite[p.~2]{Spencer2007LogicRandomStructuresAuthor}.  The same definition
applies to any sequence of probability spaces on structures of size \(n\)
defined for all sufficiently large \(n\).

A finite conjunction of properties each of which holds almost surely holds
almost surely, because \(\Pr[\neg(A_1\wedge\cdots\wedge A_s)]\leq
\sum_{i}\Pr[\neg A_i]\) and each term tends to \(0\).
\end{hbdefinition}

\begin{hbconvention}{def:p03-w03-p03-065-asymptotic-order}{The relations \(\ll\) and \(\gg\)}
For functions \(f\geq0\) and \(g>0\) of \(n\), \(f\ll g\) means
\(f(n)/g(n)\to0\) as \(n\to\infty\); and \(f\gg g\) means that \(f(n)>0\) for
all sufficiently large \(n\) and \(g\ll f\).  These
are the comparisons used in the source's definition of a threshold
\cite[pp.~2 and 4]{Spencer2007LogicRandomStructuresAuthor}.  If \(p\ll p_0\) and \(c>0\) is
constant then \(p\ll cp_0\); and \(p\ll p_0\) holds exactly when
\(p^{\,\gamma}\ll p_0^{\,\gamma}\) for any fixed \(\gamma>0\).
\end{hbconvention}

\begin{hbdefinition}{def:p03-w03-p03-065-threshold}{Threshold function}
A function \(p_0(n)\) is a \emph{threshold function} for a property \(A\) if
\(A\) holds almost never whenever \(p(n)\ll p_0(n)\), and \(A\) holds almost
surely whenever \(p(n)\gg p_0(n)\) \cite[pp.~2 and 4]{Spencer2007LogicRandomStructuresAuthor}.
The notion is due to Erd\H{o}s and R\'enyi, and the subject begins with their
1960 paper on the evolution of the random graph, as the source records
\cite[p.~4]{Spencer2007LogicRandomStructuresAuthor}.

A threshold function is not unique.  By \cref{def:p03-w03-p03-065-asymptotic-order},
if \(p_0\) is a threshold function then so is \(cp_0\) for every constant
\(c>0\); the source makes this point with \(5/n\) in place of \(n^{-1}\)
\cite[p.~2]{Spencer2007LogicRandomStructuresAuthor}.
\end{hbdefinition}

\begin{hbwarning}{def:p03-w03-p03-065-threshold-limits}{What a threshold does not provide}
The definition says nothing about \(p(n)\) of the same order as \(p_0(n)\),
where the limiting probability may lie strictly between \(0\) and \(1\) or
fail to exist.  A property need not be monotone in \(p\), a threshold function
need not exist, and the probabilities need not converge at all
\cite[pp.~2--3]{Spencer2007LogicRandomStructuresAuthor}.  For example, at \(p=n^{-1/2}\) the
property ``every two vertices have a common neighbour'' holds almost never
(\cref{thm:p03-w03-p03-065-common-neighbour}), although a fixed pair of
vertices has, in expectation, about one common neighbour.
\end{hbwarning}

The threshold examples below are proved by counting.  The two counting tools
are the first and second moment bounds.

\begin{hblemma}{thm:p03-w03-p03-065-moments}{proved=annex:p03-w03-p03-065-moments}{First and second moment bounds}
Let \(X\) be a random variable with values in \(\{0,1,2,\ldots\}\) and finite
second moment.  Then \(\Pr[X>0]\leq\mathbb E[X]\).  If \(\mathbb E[X]>0\),
then
\[
  \Pr[X=0]\leq\frac{\operatorname{Var}[X]}{\mathbb E[X]^2}.
\]
If \(X=\sum_{i\in I}\mathbf 1_{B_i}\) is a finite sum of indicators of events,
then
\[
  \operatorname{Var}[X]\leq\mathbb E[X]
  +\sum_{(i,j)\in D}\bigl(\Pr[B_i\wedge B_j]-\Pr[B_i]\Pr[B_j]\bigr),
\]
where \(D\) is any set of ordered pairs \((i,j)\), \(i\neq j\), containing
every pair for which \(B_i\) and \(B_j\) are not independent.
\end{hblemma}

\begin{hbtheorem}{thm:p03-w03-p03-065-clique}{proved=annex:p03-w03-p03-065-clique}{Threshold for a complete subgraph}
Let \(k\geq2\) be fixed, and let \(A_k\) be the property ``there are \(k\)
distinct pairwise adjacent vertices'', the first-order sentence
\(\exists x_1\cdots\exists x_k\bigwedge_{i<j}(x_i\neq x_j\wedge x_i\sim x_j)\).
If \(p\ll n^{-2/(k-1)}\), then \(A_k\) holds almost never.  If
\(p\gg n^{-2/(k-1)}\), then \(A_k\) holds almost surely.  Thus
\(n^{-2/(k-1)}\) is a threshold function for \(A_k\).  In particular
\(n^{-1}\) is a threshold function for containing a triangle and
\(n^{-2/3}\) is one for containing a \(K_4\)
\cite[pp.~4--5]{Spencer2007LogicRandomStructuresAuthor}.
\end{hbtheorem}

For \(K_4\) the source computes the expected number of copies,
asymptotically \(n^4p^6/24\), and observes that its tending to infinity does
not by itself give a copy; it appeals to the variance
\cite[p.~4]{Spencer2007LogicRandomStructuresAuthor}.  The annex carries out that variance
estimate for every \(k\).

\begin{hbtheorem}{thm:p03-w03-p03-065-isolated}{proved=annex:p03-w03-p03-065-isolated}{Threshold for having no isolated vertex}
Let \(I\) be the property \(\forall x\,\exists y\;x\sim y\), ``no vertex is
isolated'', and let \(\varepsilon>0\) be fixed.
\begin{enumerate}
\item If \(pn\geq(1+\varepsilon)\ln n\) for all sufficiently large \(n\),
then \(I\) holds almost surely.
\item If \(\varepsilon<1\) and \(pn\leq(1-\varepsilon)\ln n\) for all
sufficiently large \(n\), then \(I\) holds almost never.
\end{enumerate}
Consequently \(n^{-1}\ln n\) is a threshold function for \(I\)
\cite[p.~5]{Spencer2007LogicRandomStructuresAuthor}.
\end{hbtheorem}

The source records that connectivity behaves like having no isolated vertex,
and quotes the sharper Erd\H{o}s--R\'enyi limit inside the critical window
\cite[p.~5]{Spencer2007LogicRandomStructuresAuthor}.  Connectivity is not a
first-order property of graphs; it enters here as a comparison, and the next
theorem shows that at this scale it fails only through isolated vertices.

\begin{hbtheorem}{thm:p03-w03-p03-065-connectivity}{proved=annex:p03-w03-p03-065-connectivity}{Connectivity in the critical window}
Let \(c\) be a real constant and \(p(n)=(\ln n+c)/n\), for all \(n\) large
enough that \(0\leq p(n)\leq1\).  Let \(X\) be the number of isolated
vertices of \(G(n,p)\).  Then
\[
  \lim_{n\to\infty}\Pr[X=0]=e^{-e^{-c}}
  \quad\text{and}\quad
  \lim_{n\to\infty}\Pr[G(n,p)\text{ is connected}]=e^{-e^{-c}}
\]
\cite[p.~5]{Spencer2007LogicRandomStructuresAuthor}.  So inside this window
neither \(I\) nor connectivity holds almost surely or almost never.
\end{hbtheorem}

\begin{hbtheorem}{thm:p03-w03-p03-065-common-neighbour}{proved=annex:p03-w03-p03-065-common-neighbour}{Threshold for common neighbours}
Let \(C\) be the property
\(\forall x_1\forall x_2\,\bigl(x_1\neq x_2\rightarrow\exists y\,(y\sim x_1\wedge
y\sim x_2)\bigr)\), ``every two vertices have a common neighbour''.
\begin{enumerate}
\item If \(\varepsilon>0\) and \(p^2n\geq(2+\varepsilon)\ln n\) for all
sufficiently large \(n\), then \(C\) holds almost surely.
\item If \(0<\varepsilon<2\) and \(p^2n\leq(2-\varepsilon)\ln n\) for all
sufficiently large \(n\), then \(C\) holds almost never.
\end{enumerate}
Consequently \(n^{-1/2}\ln^{1/2}n\) is a threshold function for \(C\), and
\(C\) holds almost never when \(p=n^{-1/2}\)
\cite[p.~5]{Spencer2007LogicRandomStructuresAuthor}.
\end{hbtheorem}

The source explains the logarithmic factor.  A fixed pair has about \(np^2\)
common neighbours in expectation, which suggests \(n^{-1/2}\).  Making every
one of the roughly \(n^2/2\) pairs succeed costs the extra factor
\(\ln^{1/2}n\) \cite[p.~5]{Spencer2007LogicRandomStructuresAuthor}.  The source treats
``every two vertices are joined by a path of length three'' in the same way.
The expected number of such paths between a fixed pair is about \(n^2p^3\),
and it gives \(n^{-2/3}\ln^{1/3}n\) as the threshold function
\cite[p.~5]{Spencer2007LogicRandomStructuresAuthor}.  That threshold, with
sharp constants on both sides, is proved in
\cref{thm:p03-w03-p03-064-path3}.

\subsection{The sparse regime \(p=n^{-\alpha}\)}
\label{sec:p03-w03-p03-065-sparse-regime}

In each threshold above, the power of \(n\) is rational, and any further
factors are of smaller order.  Spencer and Shelah gave a formal explanation in
1988 for the absence of natural thresholds at irrational powers, in the form
of a zero--one law \cite[p.~5]{Spencer2007LogicRandomStructuresAuthor}.

\begin{hbtheorem}{thm:p03-w03-p03-065-shelah-spencer}{proved=annex:p03-w03-p03-065-shelah-spencer}{Zero--one law for \(p=n^{-\alpha}\), \(\alpha\) irrational}
Let \(0<\alpha<1\) be irrational and \(p(n)=n^{-\alpha}\).  For every
first-order sentence \(A\) in the language of graphs,
\[
  \lim_{n\to\infty}\Pr[G(n,p)\entails A]=0\quad\text{or}\quad 1
\]
\cite[pp.~5--6]{Spencer2007LogicRandomStructuresAuthor}.
\end{hbtheorem}

The source's route to this theorem is the one it uses for the unary model of
its Section~1.  It finds a theory \(T_\alpha\) each of whose members holds
almost surely and proves \(T_\alpha\) complete by countable models and the
Ehrenfeucht game \cite[pp.~3 and 6]{Spencer2007LogicRandomStructuresAuthor}.
The closure construction, the almost sure theory and the proof of
completeness are given in the source's Sections~4 and 5
\cite[pp.~8--12]{Spencer2007LogicRandomStructuresAuthor}.  The proof given
here has two halves.  The probabilistic half, that every member of
\(T_\alpha\) holds almost surely, rests on the extension statements of the
next section and is completed by the generic extension theorem
(\cref{thm:p03-w03-p03-065-generic-extension}).  The logical half, that
\(T_\alpha\) is complete, is Duplicator's closure strategy in
\cref{thm:p03-w03-p03-072-sparse-zero-one}.  The source sets aside
\(\alpha>1\) as considerably simpler, and does not consider it
\cite[p.~6]{Spencer2007LogicRandomStructuresAuthor}.  The hypothesis that
\(\alpha\) is irrational is used through the dense--sparse dichotomy of the
next section; \cref{def:p03-w03-p03-065-rational-boundary} shows what goes
wrong for one extension statement at a rational exponent.

\section{Extension statements}
\label{sec:p03-w03-p03-065-extension-statements}

The properties ``every vertex has a neighbour'', ``every two vertices have a
common neighbour'' and ``every two vertices are joined by a path of length
three'' share one form.  For all \(x_1,\ldots,x_r\) there exist
\(y_1,\ldots,y_v\) with certain prescribed adjacencies, each between two
\(y\)'s or between an \(x\) and a \(y\).  No adjacency between two \(x\)'s
is mentioned, and no nonadjacency is demanded
\cite[p.~6]{Spencer2007LogicRandomStructuresAuthor}.

\begin{hbdefinition}{def:p03-w03-p03-065-rooted-graph}{Rooted graph and extension statement}
A \emph{rooted graph} is a pair \((R,H)\) where \(H\) is a finite graph with
vertex set \(V(H)\) and edge set \(E(H)\), and \(R\) is a proper subset of
\(V(H)\).  The elements of \(R\) are the \emph{roots}, labelled
\(x_1,\ldots,x_r\), and the others are the \emph{nonroots}, labelled
\(y_1,\ldots,y_v\).  The parameters of \((R,H)\) are
\begin{itemize}
\item \(r=|R|\), the number of roots;
\item \(v=|V(H)\setminus R|\), the number of nonroots, so that \(v>0\);
\item \(e\), the number of edges of \(H\) not joining two roots
\end{itemize}
\cite[p.~6]{Spencer2007LogicRandomStructuresAuthor}.  The case \(r=0\) is allowed.

The \emph{extension statement} \(\operatorname{Ext}(R,H)\) is the
first-order sentence saying that for all distinct \(x_1,\ldots,x_r\) there
exist distinct \(y_1,\ldots,y_v\), distinct from the \(x\)'s, such that
\(y_i\sim y_j\) whenever \(y_iy_j\in E(H)\) and \(x_i\sim y_j\) whenever
\(x_iy_j\in E(H)\).  Edges between roots are not examined, and extra edges are
allowed \cite[p.~6]{Spencer2007LogicRandomStructuresAuthor}.  For a graph \(G\) and distinct
vertices \(x_1,\ldots,x_r\) of \(G\), a tuple \(y_1,\ldots,y_v\) with these
properties is an \emph{\((R,H)\) extension over} \(x_1,\ldots,x_r\).  The
source writes the quantifiers without the word ``distinct''; the reading
adopted here is the one under which \(x_1,\ldots,x_r,y_1,\ldots,y_v\) are
the distinct labelled vertices of \(H\).  For \(r=0\) the statement is purely
existential.
\end{hbdefinition}

For \(S\subseteq V(H)\) write \(H|_S\) for the restriction of \(H\) to \(S\),
the graph on \(S\) whose edges are the edges of \(H\) with both ends in \(S\).

\begin{hbdefinition}{def:p03-w03-p03-065-dense-sparse}{Dense, sparse, rigid, safe}
Fix an irrational \(\alpha\in(0,1)\), and let \((R,H)\) have parameters
\(r,v,e\).  Then \((R,H)\) is
\begin{itemize}
\item \emph{dense} if \(v-e\alpha<0\), and \emph{sparse} if \(v-e\alpha>0\);
\item \emph{rigid} if \((S,H)\) is dense for every \(S\) with
\(R\subseteq S\subsetneq V(H)\);
\item \emph{safe} if \((R,H|_S)\) is sparse for every \(S\) with
\(R\subsetneq S\subseteq V(H)\)
\end{itemize}
\cite[p.~6]{Spencer2007LogicRandomStructuresAuthor}.  The rooted graph \((R,H|_S)\) is a
\emph{subextension} of \((R,H)\), and \((S,H)\) is a \emph{nailextension}
of \((R,H)\), in which more vertices are nailed down as roots
\cite[p.~7]{Spencer2007LogicRandomStructuresAuthor}.  Taking \(S=R\) shows that rigid
implies dense, and taking \(S=V(H)\) shows that safe implies sparse.
Roughly, rigid means dense throughout and safe means sparse throughout.
\end{hbdefinition}

Since \(v>0\), the number \(v-e\alpha\) is never \(0\): if \(e=0\) it equals
\(v\), and if \(e>0\) it vanishes only when \(\alpha=v/e\) is rational.  So
for irrational \(\alpha\) every rooted graph is exactly one of dense and
sparse.  This is where irrationality enters \cite[p.~6]{Spencer2007LogicRandomStructuresAuthor}.

\begin{hbexample}{def:p03-w03-p03-065-pi-seven}{Classifications at \(\alpha=\pi/7\)}
The source fixes \(\alpha=\pi/7=0.4487\ldots\) for its examples, as a value
with no special properties \cite[p.~7]{Spencer2007LogicRandomStructuresAuthor}.  Below, a
rooted graph is listed by its extension statement.
\begin{enumerate}
\item \emph{Every two vertices have a common neighbour}: \(y_1\) adjacent to
\(x_1,x_2\); \(r=2\), \(v=1\), \(e=2\), and \(1-2\alpha>0\).  The only
subextension is the whole, so the rooted graph is sparse and safe.  The
source's wording at this point reads ``a neighbor''; the rooted graph it
describes is the common-neighbour one.
\item \emph{Every three vertices have a common neighbour}: \(r=3\), \(v=1\),
\(e=3\), \(1-3\alpha<0\); dense, and rigid because the only nailextension
with \(S\subsetneq V(H)\) is \(S=R\).
\item \emph{Every vertex lies in a \(K_5\)}: \(r=1\), \(v=4\), \(e=10\),
\(4-10\alpha<0\).  Nailing down \(4-k\) of the nonroots leaves \(k\) nonroots
and \(\binom k2+k(5-k)=k(9-k)/2\) counted edges.  For \(k=1,2,3,4\) the values
\(k-\alpha k(9-k)/2\) are \(1-4\alpha\), \(2-7\alpha\), \(3-9\alpha\) and
\(4-10\alpha\), all negative.  So the rooted graph is dense and rigid.
\item \emph{Every vertex lies in a \(K_4\)}: \(y_1,y_2,y_3,x_1\) with all six
adjacencies, so \(r=1\), \(v=3\), \(e=6\), and \(3-6\alpha=0.307\ldots>0\).
The subextensions with \(a=1,2,3\) nonroots have \(\binom a2+a\) counted
edges, giving \(1-2\alpha\), \(2-3\alpha\) and \(3-6\alpha\), all positive.
So this rooted graph is sparse and safe.  The source lists \(e=10\) here and
calls the rooted graph dense and rigid while recording \(v-e\alpha>0\)
\cite[p.~7]{Spencer2007LogicRandomStructuresAuthor}.  Those statements are inconsistent with
one another, and the correct values are the ones just computed.
\item \emph{Every two vertices lie in a \(K_4\), except that they may be
nonadjacent}: \(y_1,y_2,x_1,x_2\) with the five adjacencies other than
\(x_1x_2\); \(r=2\), \(v=2\), \(e=5\), \(2-5\alpha<0\).  The one proper
nailextension, nailing \(y_1\) or \(y_2\), has \(v=1\), \(e=3\) and
\(1-3\alpha<0\).  So the rooted graph is dense and rigid.
\item \emph{Every three vertices have a common neighbour which itself has a
different neighbour}: \(y_1\) adjacent to \(x_1,x_2,x_3\) and \(y_2\)
adjacent to \(y_1\); \(r=3\), \(v=2\), \(e=4\), \(2-4\alpha>0\), so it is
sparse.  It is not safe: the subextension on
\(S=\{x_1,x_2,x_3,y_1\}\) is the dense one of item~2.
\item \emph{Every four vertices have a common neighbour which itself has a
different neighbour}: \(r=4\), \(v=2\), \(e=5\), \(2-5\alpha<0\), so it is
dense.  It is not rigid: nailing down \(y_1\) gives \(r=5\), \(v=1\), \(e=1\)
and \(1-\alpha>0\).
\end{enumerate}
\end{hbexample}

The expected number of \((R,H)\) extensions over fixed roots is of order
\(n^vp^e=n^{v-e\alpha}\) \cite[p.~7]{Spencer2007LogicRandomStructuresAuthor}.  This makes
the next theorem plausible, but the expectation alone does not prove it:
extensions overlap, so their indicator events are not independent.  The
proof uses the following two inequalities.  Both concern finitely many
independent \(\{0,1\}\)-valued coordinates, here the indicators of the
\(\binom n2\) potential edges.  An event is \emph{increasing} if it is
preserved when coordinates are changed from \(0\) to \(1\), and
\emph{decreasing} if its complement is increasing.

\begin{hblemma}{thm:p03-w03-p03-065-harris}{proved=annex:p03-w03-p03-065-harris}{Harris inequality}
Let \(\omega=(\omega_1,\ldots,\omega_m)\) have independent coordinates with
values in \(\{0,1\}\).  If \(f,g\) are real functions of \(\omega\), \(f\)
nondecreasing and \(g\) nonincreasing in each coordinate, then
\(\mathbb E[fg]\leq\mathbb E[f]\,\mathbb E[g]\).  In particular, if \(B\) is
increasing and \(D\) is decreasing, then \(\Pr[B\wedge D]\leq\Pr[B]\Pr[D]\).
\end{hblemma}

\begin{hblemma}{thm:p03-w03-p03-065-janson}{proved=annex:p03-w03-p03-065-janson}{Janson inequality}
Let \(\omega\) be as in \cref{thm:p03-w03-p03-065-harris}, with
\(\Pr[\omega_k=1]<1\) for every \(k\).  Let \(\Omega\) be a finite index
set and, for \(i\in\Omega\), let \(E_i\) be a set of coordinates and \(B_i\)
the event that \(\omega_k=1\) for every \(k\in E_i\).  Write \(i\approx j\)
if \(i\neq j\) and \(E_i\cap E_j\neq\emptyset\), and set
\[
  \mu=\sum_{i\in\Omega}\Pr[B_i],\qquad
  \Delta=\sum_{(i,j):\,i\approx j}\Pr[B_i\wedge B_j],
\]
the second sum over ordered pairs.  Then
\[
  \Pr\Bigl[\bigwedge_{i\in\Omega}\neg B_i\Bigr]\leq e^{-\mu+\Delta/2},
\]
and, if \(\Delta\geq\mu>0\),
\[
  \Pr\Bigl[\bigwedge_{i\in\Omega}\neg B_i\Bigr]\leq e^{-\mu^2/(2\Delta)}.
\]
\end{hblemma}

\begin{hbtheorem}{thm:p03-w03-p03-065-safe-criterion}{proved=annex:p03-w03-p03-065-safe-criterion}{Almost sure extension statements}
Let \(0<\alpha<1\) be irrational, \(p=n^{-\alpha}\), and let \((R,H)\) be a
rooted graph.  Then \(\operatorname{Ext}(R,H)\) holds almost surely in
\(G(n,p)\) if \((R,H)\) is safe, and holds almost never if \((R,H)\) is not
safe \cite[p.~7]{Spencer2007LogicRandomStructuresAuthor}.  More precisely, if \((R,H)\) is
safe there are constants \(c,\varepsilon>0\) such that, for all sufficiently
large \(n\), the probability that some distinct \(x_1,\ldots,x_r\) have no
\((R,H)\) extension is at most \(n^re^{-cn^{\varepsilon}}\).
\end{hbtheorem}

The source proves the second half by expectation.  If a subextension
\((R,H|_S)\) is dense, then a fixed root tuple has an expected number of
\(H|_S\) extensions tending to \(0\), and every \((R,H)\) extension restricts
to an \((R,H|_S)\) extension \cite[p.~7]{Spencer2007LogicRandomStructuresAuthor}.  The
exponential bound in the first half is what the union bound over the
\(O(n^r)\) root tuples requires.  The source states a two-sided version of it
for the number of extensions, with failure probability \(o(n^{-r})\), as the
technical heart of its almost sure theory \cite[pp.~10--11]{Spencer2007LogicRandomStructuresAuthor}.

\begin{hbexample}{def:p03-w03-p03-065-rigid-special}{Rigid extensions as special properties}
At \(\alpha=\pi/7\), three randomly chosen vertices almost surely have no
common neighbour, and the rooted graph of item~2 of
\cref{def:p03-w03-p03-065-pi-seven} is not safe.  So by
\cref{thm:p03-w03-p03-065-safe-criterion} it is almost surely false that
\emph{every} three vertices have a common neighbour.  Yet almost surely
\emph{some} three vertices do.  The rooted graph \((\emptyset,K_{1,3})\), a
star with centre \(y_1\) and leaves \(y_2,y_3,y_4\), has \(r=0\), \(v=4\),
\(e=3\).  Every nonempty set of \(a\) of its vertices spans at most \(a-1\)
edges, because the star is a forest, so each subextension has
\(a-(a-1)\alpha>0\), and the rooted graph is safe.  By the theorem a star
with three leaves almost surely exists, and its three leaves have a common
neighbour.  Having a common neighbour is thus a special property of a triple,
which only some triples have.  By contrast, every pair of vertices has a common
neighbour, so having one is not special for a pair.  The source anticipates
that the special properties of bounded sets of vertices are all described by
rigid extensions \cite[pp.~7--8]{Spencer2007LogicRandomStructuresAuthor}.
\end{hbexample}

\begin{hbwarning}{def:p03-w03-p03-065-rational-boundary}{The criterion needs irrational \(\alpha\)}
At a rational \(\alpha\) a rooted graph can have \(v-e\alpha=0\).  It is then
neither dense nor sparse, and \cref{thm:p03-w03-p03-065-safe-criterion} does
not apply.  For the common-neighbour rooted graph and \(\alpha=1/2\), the
expected number of extensions over a fixed pair is \((n-2)n^{-1}\to1\).  It
does not tend to \(0\), and it does not grow like a positive power of \(n\).
The statement nevertheless holds almost never there, by
\cref{thm:p03-w03-p03-065-common-neighbour}.  The behaviour at such
\(\alpha\) depends on lower-order factors that the dense--sparse
classification does not see.
\end{hbwarning}

\section{Counting extensions and generic extensions}
\label{sec:p03-w03-p03-065-counting}

The almost sure theory of \(G(n,n^{-\alpha})\) asks more of an extension
than that it exists.  The new vertices should have no further adjacencies
than the rooted graph demands, and no special properties relative to the
roots beyond those the rooted graph forces.  The source expresses the second
demand by rigid extensions, through the notion of a \(t\)-generic extension,
and derives it from a large deviation estimate for the number of extensions,
whose proof it omits \cite[pp.~10--11]{Spencer2007LogicRandomStructuresAuthor}.
This section proves the estimate in the form the source uses, with failure
probability \(o(n^{-r})\) for each tuple of \(r\) roots, and derives the
generic extension statement from it.  Throughout the section
\(\alpha\in(0,1)\) is irrational and \(p=n^{-\alpha}\).

\begin{hbdefinition}{def:p03-w03-p03-065-exponent}{The exponent of a rooted graph; extension counts}
Let \(F\) be a finite graph and \(T\subseteq V(F)\).  Write \(e_F(T)\) for
the number of edges of \(F\) with at least one endpoint outside \(T\), and
put
\[
  \lambda_F(T)=|V(F)\setminus T|-\alpha\,e_F(T).
\]
If \(T\subsetneq V(F)\), then \((T,F)\) is a rooted graph with
\(v=|V(F)\setminus T|\) and \(e=e_F(T)\), so \(\lambda_F(T)=v-e\alpha\); and
\(\lambda_F(V(F))=0\).  In this notation a rooted graph \((R,F)\) is rigid
when \(\lambda_F(T)<0\) for every \(T\) with \(R\subseteq T\subsetneq V(F)\),
and safe when \(\lambda_{F|_S}(R)>0\) for every \(S\) with
\(R\subsetneq S\subseteq V(F)\) (\cref{def:p03-w03-p03-065-dense-sparse}).

For a graph \(G\), a rooted graph \((R,F)\) with \(r\) roots, and distinct
vertices \(x=(x_1,\ldots,x_r)\) of \(G\), let \(N_{R,F}(x)\) be the number of
\((R,F)\) extensions over \(x\), counted as tuples.  In \(G(n,p)\) it is a
random variable.  If \((R,F)\) has \(w\) nonroots and \(f\) counted edges,
write
\[
  \mu_{R,F}=(n-r)(n-r-1)\cdots(n-r-w+1)\,p^{f}.
\]
\end{hbdefinition}

\begin{hblemma}{thm:p03-w03-p03-065-decomposition}{proved=annex:p03-w03-p03-065-decomposition}{Exponent calculus and the rigid--safe decomposition}
Let \((R,F)\) be a rooted graph.
\begin{enumerate}
\item If \(T\subseteq T'\subseteq V(F)\), then
  \(\lambda_F(T)=\lambda_{F|_{T'}}(T)+\lambda_F(T')\).
\item If \(R\subsetneq S\subsetneq S'\subseteq V(F)\), \((R,F|_S)\) is rigid
  and \((S,F|_{S'})\) is rigid, then \((R,F|_{S'})\) is rigid.
\item There is a set \(S\) with \(R\subseteq S\subseteq V(F)\) such that
  \((R,F|_S)\) is rigid unless \(S=R\), and \((S,F)\) is safe unless
  \(S=V(F)\).
\end{enumerate}
\end{hblemma}

\begin{hbtheorem}{thm:p03-w03-p03-065-counting}{proved=annex:p03-w03-p03-065-counting}{Counting theorem for safe extensions}
Let \((R,H)\) be safe with parameters \(r,v,e\), and \(\mu=\mu_{R,H}\).
Then \(\mathbb E[N_{R,H}(x)]=\mu\) for every tuple \(x\) of \(r\) distinct
vertices, and \(2^{-v}n^{v-e\alpha}\leq\mu\leq n^{v-e\alpha}\) when
\(n\geq2(r+v)\).  For every \(\varepsilon>0\) and every fixed tuple \(x\) of
distinct vertices,
\[
  \Pr\bigl[|N_{R,H}(x)-\mu|>\varepsilon\mu\bigr]=o(n^{-r}),
\]
the bound being uniform in \(x\).  Consequently, almost surely,
\((1-\varepsilon)\mu\leq N_{R,H}(x)\leq(1+\varepsilon)\mu\) for all tuples
\(x\) of \(r\) distinct vertices at once
\cite[pp.~10--11]{Spencer2007LogicRandomStructuresAuthor}.
\end{hbtheorem}

The source calls the displayed estimate a large deviation result and the
consequence a Counting Theorem.  It also remarks that the failure probability
can be bounded by an exponentially small quantity
\cite[p.~10]{Spencer2007LogicRandomStructuresAuthor}.  That stronger bound is
proved here only for the existence of one extension
(\cref{thm:p03-w03-p03-065-safe-criterion}); for the count, the
\(o(n^{-r})\) form is proved, which is the form the union bound over root
tuples uses.

\begin{hblemma}{thm:p03-w03-p03-065-rigid-count}{proved=annex:p03-w03-p03-065-rigid-count}{Rigid extensions are boundedly many}
Let \((R,F)\) be rigid with \(r\) roots.  There is a constant \(K\) such
that, almost surely, \(N_{R,F}(x)\leq K\) for every tuple \(x\) of \(r\)
distinct vertices.
\end{hblemma}

\begin{hbcorollary}{thm:p03-w03-p03-065-upper-count}{proved=annex:p03-w03-p03-065-upper-count}{An upper bound for every rooted graph}
Let \((R,F)\) be a rooted graph with \(r\) roots and let \(S\) be a set as
in item~3 of \cref{thm:p03-w03-p03-065-decomposition}.  There is a constant
\(C\) such that, almost surely, \(N_{R,F}(x)\leq C\,n^{\lambda_F(S)}\) for
every tuple \(x\) of \(r\) distinct vertices.
\end{hbcorollary}

The generic extension statement uses the notion of a \(t\)-generic tuple and
the Generic Extension schema of \cref{def:p03-w03-p03-070-theory}.  For a
safe \((R,H)\) and \(t\geq0\), the schema asks that every tuple of distinct
roots have an \((R,H)\) extension with (1) the edges of \(H\), (2) no further
edges among its nonroots or between a root and a nonroot, and (3) for
\(t\geq1\), no set of at most \(t\) further vertices that is rigid over roots
and nonroots together and has a vertex adjacent to a nonroot
\cite[p.~10]{Spencer2007LogicRandomStructuresAuthor}.

\begin{hbtheorem}{thm:p03-w03-p03-065-generic-extension}{proved=annex:p03-w03-p03-065-generic-extension}{Generic extension}
Let \((R,H)\) be safe with parameters \(r,v,e\), and let \(t\geq0\).  Almost
surely in \(G(n,n^{-\alpha})\), every tuple of \(r\) distinct vertices has an
\((R,H)\) extension satisfying (1), (2) and (3) of the Generic Extension
schema of \cref{def:p03-w03-p03-070-theory}.  More precisely, there are
constants \(C,\gamma>0\) such that, almost surely, for every such tuple \(x\)
the number of \((R,H)\) extensions over \(x\) that violate (2) or (3) is at
most \(Cn^{v-e\alpha-\gamma}\), while the number of all \((R,H)\) extensions
over \(x\) is at least \(\mu_{R,H}/2\geq2^{-v-1}n^{v-e\alpha}\)
\cite[pp.~10--11]{Spencer2007LogicRandomStructuresAuthor}.
\end{hbtheorem}

\begin{hbexample}{def:p03-w03-p03-065-generic-example}{Generic common neighbours at \(\alpha=\pi/7\)}
Take the common-neighbour rooted graph \((R,H)\), \(y_1\) adjacent to
\(x_1,x_2\), at \(\alpha=\pi/7=0.4487\ldots\), so that
\(2\alpha=0.897\ldots<1<3\alpha=1.346\ldots\).  It is safe (item~1 of
\cref{def:p03-w03-p03-065-pi-seven}), and by
\cref{thm:p03-w03-p03-065-counting} almost surely every two vertices have
\((1+o(1))\,n^{1-2\alpha}\) common neighbours, where
\(1-2\alpha=0.102\ldots\).  Condition (2) is empty here: the only pairs
involving \(y_1\) and a root are the two edges of \(H\).  For \(t=1\), a
single vertex \(z\) is rigid over \(\{x_1,x_2,y_1\}\) when its number \(k\)
of neighbours there satisfies \(1-k\alpha<0\), that is \(k=3\).  So \(y_1\)
fails to be \(1\)-generic over \(x_1,x_2\) exactly when some \(z\) is
adjacent to \(x_1\), \(x_2\) and \(y_1\).  The rooted graph
\((\{x_1,x_2\},F)\), with \(F\) the edges \(y_1x_1,y_1x_2,zx_1,zx_2,zy_1\),
has \(\lambda_F(\{x_1,x_2\})=2-5\alpha<0\), and each of its two
nailextensions with one root added has exponent \(1-3\alpha<0\), so it is
rigid.  By \cref{thm:p03-w03-p03-065-rigid-count}, almost surely every pair
has at most a constant number of common neighbours that are not
\(1\)-generic, against about \(n^{0.102}\) common neighbours in all.
\end{hbexample}

\begin{hbwarning}{def:p03-w03-p03-065-generic-limits}{What the counting and generic extension theorems do not give}
The constants \(C\), \(\gamma\) and the size of \(n\) beyond which the bounds
hold depend on \((R,H)\), \(t\) and \(\alpha\).  The theorems give no rate
uniform over all safe rooted graphs; each sentence of the schema concerns one
rooted graph and one \(t\), so none is used.  The counting theorem needs
safety, not only sparseness.  For item~6 of
\cref{def:p03-w03-p03-065-pi-seven}, which is sparse but not safe,
\(\mu\) tends to infinity, yet a fixed triple of roots almost never has an
extension at all, because the dense subextension of item~2 already fails
(\cref{thm:p03-w03-p03-065-safe-criterion}).
\end{hbwarning}

\section{Relationships}
\label{sec:p03-w03-p03-065-relationships}

For constant \(p\in(0,1)\) the corresponding computation is simpler and is
the one in \cref{sec:p03-zero-one-law}.  One new vertex with prescribed
adjacencies to \(r\) given vertices \emph{and} prescribed nonadjacencies to
\(s\) others is a witness with probability \(p^r(1-p)^s\), independently for
different candidates, and the failure bound
\(n^{r+s}(1-p^r(1-p)^s)^{n-r-s}\) tends to \(0\)
\cite[pp.~12--13]{Spencer2007LogicRandomStructuresAuthor}.  Those extension axioms
axiomatize the theory of the countable random graph
(\cref{thm:p03-random-graph}).  The labelled zero--one law and its rate for
\(p=1/2\) are proved from Fagin's paper in
\cref{thm:p09-fagin1976-labelled}.  The completeness arguments by games use
the Ehrenfeucht--Fra\"iss\'e theorem, \cref{thm:p03-ef}.  The countable
models of the almost sure theories are classified as follows.  For constant
\(p\) the almost sure theory is \(\aleph_0\)-categorical, its countable model
being the Rado graph (\cref{thm:p03-w03-p03-067-categorical}).  For
irrational \(\alpha\) the theory \(T_\alpha\), whose members hold almost
surely by \cref{thm:p03-w03-p03-065-generic-extension}, is not
\(\aleph_0\)-categorical (\cref{thm:p03-w03-p03-067-sparse-models}); the
estimate that the argument there takes from the source is
\cref{thm:p03-w03-p03-065-counting}.  The theory \(T_\alpha\) and the closure
lemmas used by Duplicator are in \cref{def:p03-w03-p03-070-theory},
\cref{thm:p03-w03-p03-070-weights} and
\cref{thm:p03-w03-p03-070-finite-closure}.

Extension statements differ from those extension axioms in two ways.  They
never demand a nonadjacency, and their almost sure truth depends on the
exponent \(\alpha\) through the sign of \(a-e_a\alpha\) for each
subextension with \(a\) nonroots and \(e_a\) counted edges.  For a fixed
rooted graph there are finitely many such pairs \((a,e_a)\), and each sign
can change only at the rational number \(a/e_a\).  So whether a fixed
extension statement holds almost surely is constant on each open interval of
irrational \(\alpha\) avoiding these finitely many rationals, and can change
only where \(\alpha\) crosses one of them.

\section{Sources}
\label{sec:p03-w03-p03-065-sources}

All definitions, examples and source claims above are taken from
\cite{Spencer2007LogicRandomStructuresAuthor}, Sections~1--3 and 5, with the
page locators given at each use.  The Harris and Janson inequalities, the
counting theorem and all proofs in the annex are written for this chapter and
do not follow a source text.  The source states the counting theorem without
proof and outlines the proof of the generic extension statement.
